% Part: second-order-logic
% Chapter: metatheory
% Section: undecidability-and-axiomatizability

\documentclass[../../../include/open-logic-section]{subfiles}

\begin{document}

\olfileid{sol}{met}{nax}

\olsection{منطق الرتبة الثانية غير قابل للمحورة}

\begin{thm}
\ollabel{thm:sol-undecidable}
منطق الرتبة الثانية غير قابل للقرار.
\end{thm}

\begin{proof}
تصدق~!!{sentence} من الرتبة الأولى منطقيًا في منطق الرتبة الأولى إذا
وفقط إذا صدقت منطقيًا في منطق الرتبة الثانية، ومنطق الرتبة الأولى غير
قابل للقرار.
\end{proof}

\begin{thm}
\ollabel{cor:sol-not-axiomatizable}
لا يوجد لمنطق الرتبة الثانية نظام~!!{derivation} سليم وكامل.
\end{thm}

\begin{proof}
لتكن~$!A$ !!a{sentence} في لغة الحساب. تكون~$\Sat{N}{!A}$ إذا وفقط
إذا كانت~$\Th{PA^2} \Entails !A$. ولتكن~$!P$ اقتران البديهيات التسع
لـ$\Th{PA^2}$. تكون~$\Th{PA^2} \Entails !A$ إذا وفقط إذا كانت
$\Entails !P \lif !A$؛ أي~$\Sat{M}{!P \lif !A}$. تأمل الآن
الـ!!{sentence}
$\lforall[z][\lforall[u][\lforall[u'][\lforall[u''][\lforall[L][(!P' \lif !A')]]]]]$
الناتجة من استبدال~$z$ بـ$\Obj{0}$، ومتغير الدالة الأحادي الموضع~$u$
بـ$\prime$، ومتغيري الدالة الثنائيي الموضع~$u'$ و$u''$ بـ$+$ و$\times$
على الترتيب، ومتغير العلاقة الثنائي الموضع~$L$ بـ$<$، ثم التكميم الكلي
عليها. وتصدق هذه الجملة من منطق الرتبة الثانية الخالص منطقيًا إذا وفقط
إذا صدقت الجملة الأصلية منطقيًا، وإذا وفقط إذا كانت
$\Th{PA^2} \Entails !A$، وإذا وفقط إذا كانت~$\Sat{N}{!A}$. ومن ثم، لو
وجد نظام برهان سليم وكامل لمنطق الرتبة الثانية، لاستطعنا استعماله
لتعريف تعداد قابل للحساب~$f\colon \Nat \to \Sent[L_A]$
للـ!!{sentence}s الصادقة في~$\Struct{N}$. ولكانت هذه الدالة قابلة
للتمثيل في~$\Th{Q}$ بصيغة من الرتبة الأولى~$!B_f(x, y)$. وعندئذ
لعرّفت الـ!!{formula}~$\lexists[x][!B_f(x, y)]$ مجموعة
الـ!!{sentence}s الصادقة من الرتبة الأولى في~$\Struct{N}$، مناقضةً
مبرهنة تارسكي.
\end{proof}

\end{document}
